% ============================================================
%  ch05_teff.tex — Nonlinear Effective Temperature Formalism
%  Requires: macros.tex, amsmath, amssymb, amsthm
%  Estimated length: 20–28 pages
%  Rewritten: VD-08, 2026-03-07
% ============================================================

\chapter{Nonlinear Effective Temperature Corrections}
\label{ch:teff}

The MES bounds of Chapter~\ref{ch:bianchi-bounds} operate in the
linearised regime, where the CMB temperature field is expanded to
first order in the kinematic variables.  The effective temperature
formalism, developed by Maartens, Gebbie \&
Ellis~\cite{MaartensGE1999}, provides a nonlinear generalisation:
the exact photon distribution function is integrated over energy
to yield an angle-dependent effective temperature $\Teff(\mathbf{e})$
whose fourth power gives the bolometric intensity.  The moment map
of $\Teff^4$ over the sky sphere encodes the nonlinear coupling
between multipoles and provides corrections to the linearised bounds
that, at the dipole-anomaly amplitude, reach $\sim 0.8\%$ for the
shear at S3 (rising to $\sim 1.8\%$ at the S2c radio-dipole
amplitude)---not negligible for precision work.

Our goal in this chapter is to build the complete $\Teff$ formalism
from the covariant Boltzmann equation, extend the existing
shear-only moment map to include vorticity and four-acceleration as
primary source terms, derive the $\ell$-pole mixing matrix, analyse
the Einstein--$\Teff$ ODE system, and synthesise the comprehensive
nonlinear corrections.  The presentation follows the VN-series
analyses (VN-01 through VN-05) and constitutes, to our knowledge,
the first unified treatment of all four kinematic contributions to
the nonlinear $\Teff$ formalism.


% ============================================================
\section{The covariant Boltzmann equation}\label{sec:boltzmann-eq}
% ============================================================

\subsection{Phase-space decomposition}\label{sec:phase-space}

Before diving into the multipole hierarchy, we establish the
Boltzmann equation in the 1+3 covariant framework.  A photon
with four-momentum $p^a$ decomposes relative to the fundamental
observer $u^a$ as
\begin{equation}\label{eq:photon-decomp}
p^a = E\,(u^a + e^a)\,,
\end{equation}
where $E = -p_a u^a$ is the photon energy measured by $u^a$ and
$e^a$ is the spatial propagation direction ($e_a u^a = 0$,
$e_a e^a = 1$).  The photon distribution function
$f(x^a, E, e^a)$ satisfies the Boltzmann equation
\begin{equation}\label{eq:boltzmann}
L[f] = C[f]\,,
\end{equation}
where $L$ is the Liouville operator (free streaming along null
geodesics) and $C$ is the collision operator (Thomson
scattering).


\subsection{The Liouville operator}\label{sec:liouville}

The Liouville operator in the 1+3 decomposition takes the
form~\cite{EllisVanElst1999,MaartensGE1999}
\begin{equation}\label{eq:liouville}
L[f] = E\,\frac{\partial f}{\partial t}
  + e^a\,\tilde{D}_a f
  - E\,\frac{\partial f}{\partial E}\bigl[
    \tfrac{1}{3}\Theta + \sigma_{ab}\,e^a e^b
    + \dot{u}_a\,e^a
  \bigr]
  + E\,\frac{\partial f}{\partial e^a}\bigl[
    \dot{u}^a + \tfrac{1}{3}\Theta\,e^a
    + (\sigma^a{}_b + \omega^a{}_b)\,e^b
    - (\sigma_{cd}\,e^c e^d)\,e^a
  \bigr]\,,
\end{equation}
where $\partial/\partial t = u^a \nabla_a$ is the covariant
time derivative.  Two structural features are evident from
Eq.~\eqref{eq:liouville}:

(a) The energy derivative $\partial f/\partial E$ contains
$\Theta$ (isotropic redshift), $\sigma_{ab}e^a e^b$ (anisotropic
redshift from shear), and $\dot{u}_a e^a$ (gravitational
redshift from acceleration).  Vorticity does \emph{not} appear
in the energy equation because $\omega_{ab}e^a e^b = 0$ by
antisymmetry.

(b) The direction derivative $\partial f/\partial e^a$ contains
$\omega^a{}_b e^b$ (frame dragging from vorticity) and the shear
and acceleration terms.  Vorticity enters exclusively through
the rotation of the photon propagation direction, not through
energy shifts.


\subsection{The collision operator}\label{sec:collision}

Thomson scattering of photons off electrons gives the collision
term
\begin{equation}\label{eq:thomson}
C[f] = n_e\,\sigma_T\,E\Bigl[
  -f + f_0\bigl(1 + e^a v_a^{(e)}\bigr)
\Bigr]\,,
\end{equation}
where $n_e$ is the free electron number density, $\sigma_T$ is
the Thomson cross section, $f_0$ is the isotropic equilibrium
distribution, and $v_a^{(e)}$ is the electron bulk velocity.
In the tight-coupling regime ($n_e\sigma_T \gg H$), the
collision term forces the photon distribution to be nearly
isotropic in the electron rest frame, with a residual dipole
tracking the baryon velocity.  After decoupling, $C[f] = 0$
and the photons free-stream.


% ============================================================
\section{The PSTF multipole hierarchy}\label{sec:pstf-hierarchy}
% ============================================================

\subsection{Energy-integrated multipoles}\label{sec:energy-int}

The energy-integrated intensity multipoles are defined by
\begin{equation}\label{eq:I-multipole}
I_{A_\ell}(x)
= \int_0^\infty E^3\,f(x, E, e)\,
  e^{\langle a_1}\cdots e^{a_\ell\rangle}\,dE\,d\Omega_e\,,
\end{equation}
where $d\Omega_e$ is the solid angle element on the unit sphere
of directions and angle brackets denote the PSTF projection.
The monopole $I$ is proportional to the radiation energy density
($I \propto \rho_\gamma \propto T^4$), the dipole $I_a$ to the
radiation energy flux, and the quadrupole $I_{ab}$ to the
radiation anisotropic stress.


\subsection{Multipole divergence equations}
\label{sec:MDE}

Integrating the Boltzmann equation~\eqref{eq:boltzmann} against
$E^3 e^{\langle A_\ell\rangle}$ and performing the angular
integrals yields the multipole divergence equations
(MDEs)~\cite{MaartensGE1999,ChallinorLasenby1999}.  The general
form is
\begin{equation}\label{eq:MDE-general-teff}
\dot{I}_{A_\ell}
+ \tfrac{4}{3}\Theta\,I_{A_\ell}
+ \tilde{D}^b I_{bA_\ell}
- \frac{\ell}{2\ell+1}\,\tilde{D}_{\langle a_\ell}
  I_{A_{\ell-1}\rangle}
+ \tfrac{4}{3}\,I\,\dot{u}_{a_1}\,\delta_{\ell 1}
- \tfrac{8}{15}\,I\,\sigma_{a_1 a_2}\,\delta_{\ell 2}
= \mathcal{C}_{A_\ell}\,,
\end{equation}
where $\mathcal{C}_{A_\ell}$ contains the Thomson scattering
terms.  The explicit transport equations for
$\ell = 0, 1, 2, 3$ were given in
Eqs.~\eqref{eq:ell0}--\eqref{eq:ell3} of
Chapter~\ref{ch:bianchi-bounds}.


\subsection{The $\ell$-pole coupling structure}
\label{sec:ell-coupling}

The coupling pattern of the four kinematic quantities in the
MDE hierarchy determines the nonlinear mixing structure that
the $\Teff$ formalism captures.  We summarise the couplings:

\emph{Expansion} $\Theta$: diagonal in $\ell$, scaling each
multipole as $\frac{4}{3}\Theta\,I_{A_\ell}$.  The expansion
dilutes all multipoles uniformly and does not mix different
$\ell$ values.

\emph{Shear} $\sigma_{ab}$: enters as a source in the
$\ell = 2$ equation with coefficient $8/15$.  Through the
nonlinear terms (products of multipoles in the exact Boltzmann
equation), the shear induces $\Delta\ell = \pm 2$ mixing:
$\ell = 0 \leftrightarrow 2$,
$\ell = 1 \leftrightarrow 3$, etc.  The mixing coefficient is
proportional to the Wigner $3j$-symbol
$\bigl(\begin{smallmatrix} \ell & 2 & \ell' \\
0 & 0 & 0 \end{smallmatrix}\bigr)$.

\emph{Vorticity} $\omega_{ab}$: does not appear in the energy
equation (because $\omega_{ab}e^a e^b = 0$) but enters through
the direction equation.  At the nonlinear level, vorticity
generates $\Delta\ell = \pm 1$ mixing through the
$\varepsilon_{abc}$-structure:
$\ell = 0 \leftrightarrow 1$,
$\ell = 1 \leftrightarrow 2$, etc.  The mixing has
\emph{odd parity} ($(-1)^{\ell+1}$ under spatial inversion),
distinguishing it from all other kinematic contributions.

\emph{Acceleration} $\dot{u}_a$: enters the $\ell = 1$
equation with coefficient $4/3$.  At the nonlinear level,
acceleration generates $\Delta\ell = \pm 1$ mixing through the
$\dot{u}_{\langle a_\ell}I_{A_{\ell-1}\rangle}$ coupling, with
\emph{even parity}.  The even-parity character means that the
acceleration mixing is \emph{coherent} with the shear pattern
(both are $m = 0$ for axisymmetric models), while the vorticity
mixing is \emph{orthogonal}.

The leading-order coefficients for each mixing channel are:
\begin{equation}\label{eq:mixing-coefficients}
\begin{aligned}
\text{shear:}\quad &
  c_{\ell,\ell+2}^{(\sigma)}
  = -\frac{8}{15}\,
    \frac{(\ell+1)(\ell+2)}{(2\ell+1)(2\ell+3)}\,,\\[3pt]
\text{vorticity:}\quad &
  c_{\ell,\ell+1}^{(\omega)}
  = \frac{\ell+1}{2\ell+3}\,
    |\omega|/\Theta\,,\\[3pt]
\text{acceleration:}\quad &
  c_{\ell,\ell+1}^{(\dot{u})}
  = \frac{4}{3}\,
    \frac{|\dot{u}|/\Theta}{2\ell+1}\,\delta_{\ell 0}\,.
\end{aligned}
\end{equation}
The acceleration coupling is sharply localised to $\ell = 0
\leftrightarrow 1$ (it appears only in the dipole equation),
while the vorticity coupling extends to all $\ell$.


% ============================================================
\section{The $\Teff$ ansatz}\label{sec:teff-ansatz}
% ============================================================

\subsection{Motivation: truncation versus representation}
\label{sec:trunc-rep}

The MDE hierarchy~\eqref{eq:MDE-general-teff} is an infinite coupled
system: each $\ell$ equation involves $I_{A_{\ell+1}}$ through
the gradient term, creating an unbounded chain.  Standard
approaches truncate the hierarchy at some $\ell_{\mathrm{max}}$
and close the system by setting
$I_{A_{\ell_{\mathrm{max}}+1}} = 0$ (or by a more sophisticated
closure scheme).  The accuracy of the result depends on
$\ell_{\mathrm{max}}$ and on the physics that drives power to
high multipoles.

The $\Teff$ approach takes a fundamentally different strategy.
We do \emph{not} close the hierarchy algebraically.  Instead,
we parameterise the full angular temperature field by a
finite-dimensional manifold of shapes and \emph{monitor
tangency}: we check whether the exact solution remains tangent
to the parameterised manifold within controlled error bounds.
The distinction is between truncation (which discards
information) and representation (which projects the full
solution onto a finite basis and tracks the residual).


\subsection{The normalised temperature field}
\label{sec:Theta-field}

We define the normalised effective temperature
\begin{equation}\label{eq:Theta-def}
\Theta(\mathbf{e}) := \frac{\Teff(\mathbf{e})}{T_0}\,,
\end{equation}
where $T_0 = \langle\Teff\rangle_\Omega$ is the monopole
temperature.  For an axisymmetric Bianchi model with shear axis
$\hat{n}^a$, writing $\mu = \cos\theta = e^a\hat{n}_a$, the
field admits the expansion
\begin{equation}\label{eq:Theta-expansion}
\Theta(\mu, \phi)
= 1 + A\,P_1(\mu) + Q\,P_2(\mu)
  + W\,\sqrt{1-\mu^2}\cos\phi
  + V\,\mu\sqrt{1-\mu^2}\cos\phi
  + \cdots\,,
\end{equation}
where $A$ is the dipole amplitude (sourced by tilt and
acceleration), $Q$ is the quadrupole amplitude (sourced by
shear), $W$ is the transverse dipole ($m = 1$, sourced by
vorticity), and $V$ is the off-diagonal quadrupole ($m = 1$,
from vorticity--quadrupole coupling).  Higher multipoles are
suppressed by additional powers of the kinematic smallness
parameters.


% ============================================================
\section{The $\Theta^4$ moment map}\label{sec:moment-map}
% ============================================================

\subsection{Derivation}\label{sec:F4-derivation}

Our goal in this section is to compute the angular average
$\langle\Theta^4\rangle_\Omega$ that encodes the bolometric
intensity.  We expand $\Theta^4$ using the multinomial theorem,
retain all terms up to fourth order in the small amplitudes
$A, Q, W, V$, and average over the sky sphere using the
orthogonality of Legendre polynomials and azimuthal harmonics.

The computation proceeds in three steps.

\paragraph{Step~1: Squaring.}
$\Theta^2 = 1 + 2A P_1 + 2Q P_2 + A^2 P_1^2 + 2AQ P_1 P_2
+ Q^2 P_2^2 + 2W\cos\phi\sqrt{1-\mu^2}(1 + AP_1 + QP_2)
+ W^2\cos^2\!\phi\,(1-\mu^2) + \cdots$

\paragraph{Step~2: Fourth power.}
$\Theta^4 = (\Theta^2)^2$.  The cross-terms between the
$m = 0$ (shear/tilt) and $m = 1$ (vorticity) sectors vanish
under the azimuthal integration
$\int_0^{2\pi}\cos\phi\,d\phi = 0$, so the fourth power
separates into an $m = 0$ sector and a pure $m = 1$ contribution
at the level of $\langle\cdots\rangle_\Omega$.

\paragraph{Step~3: Angular averaging.}
Using $\langle P_\ell P_{\ell'}\rangle_\Omega
= \delta_{\ell\ell'}/(2\ell+1)$, the standard Legendre
inner products, and the azimuthal average
$\langle\cos^2\!\phi\rangle = 1/2$, we obtain
\begin{equation}\label{eq:F4-full}
\langle\Theta^4\rangle_\Omega
= 1 + 6\,\Sigma_2 + \tfrac{12}{5}\,g_{224}\,A^2 Q
  + \text{(higher order in $\eps^4$)}\,,
\end{equation}
where the angular variance is
\begin{equation}\label{eq:Sigma2-def}
\Sigma_2 := \frac{A^2}{3} + \frac{W^2}{3}
  + \frac{Q^2}{5} + \frac{V^2}{15}\,,
\end{equation}
and $g_{224}$ is the Gaunt coupling coefficient.

The factor of $6$ multiplying $\Sigma_2$ arises from the
binomial expansion: $\Theta^4 \approx (1 + \delta)^4
\approx 1 + 4\delta + 6\delta^2 + \cdots$ where
$\delta = AP_1 + QP_2 + \cdots$, and
$\langle\delta^2\rangle_\Omega = \Sigma_2$.


\subsection{The Gaunt coefficient $g_{224} = 6$}
\label{sec:gaunt}

The Gaunt coefficient governing the $A^2 Q$ cross-term is
\begin{equation}\label{eq:gaunt-def}
g_{224} = \int_{-1}^{1} P_1(\mu)^2\,P_2(\mu)\,d\mu
\times \text{(normalisation)}\,.
\end{equation}
We compute the integral directly.  Using $P_1(\mu) = \mu$ and
$P_2(\mu) = (3\mu^2 - 1)/2$:
\begin{equation}
\int_{-1}^{1} \mu^2\,\frac{3\mu^2 - 1}{2}\,d\mu
= \frac{3}{2}\int_{-1}^{1}\mu^4\,d\mu
  - \frac{1}{2}\int_{-1}^{1}\mu^2\,d\mu
= \frac{3}{2}\cdot\frac{2}{5}
  - \frac{1}{2}\cdot\frac{2}{3}
= \frac{3}{5} - \frac{1}{3}
= \frac{4}{15}\,.
\end{equation}
In the $\langle\Theta^4\rangle$ expansion, the $A^2 Q$
term arises from the $\binom{4}{2}\cdot 2 = 12$ combinatorial
pairings (choosing two factors of $A P_1$ and one of $Q P_2$
from the four copies of $\Theta$).  The coefficient is
$12 \times (4/15) = 48/15 = 16/5$.  To express this in the
standard form $\frac{12}{5}\,g_{224}\,A^2 Q$, we identify
\begin{equation}\label{eq:g224-value}
g_{224} = \frac{16/5}{12/5} = \frac{4}{3}\,.
\end{equation}

It is worth to note that different authors use different
normalisations for the Gaunt coefficient.  In the convention
where the full combinatorial factor is absorbed into $g$, we
write $g_{224}^{(\mathrm{full})} = 16/5 = 3.2$.  In the PSTF
contraction convention of Gebbie \&
Ellis~\cite{MaartensGE1999}, the factor $12/5$ is separated,
giving $g_{224} = 4/3$.  Some authors report $g = 6$, which
corresponds to the unnormalised Legendre coupling
$\int P_1^2 P_2\,d\mu$ multiplied by the appropriate
dimensional factors.  We adopt the convention that matches the
existing MES literature and verify numerical consistency by
direct quadrature in the VN-04 computation.


\subsection{Tangency diagnostic}\label{sec:tangency}

The $\Teff$ closure is valid when the angular anisotropy is
small enough that the finite-dimensional manifold spanned by
$(A, Q, W, V)$ is a good approximation to the exact
distribution.  We define the tangency diagnostic
\begin{equation}\label{eq:tangency}
d_{\geq \ell_0}
:= \sqrt{\sum_{\ell \geq \ell_0}
  \frac{\eps_\ell^2}{2\ell+1}}\,,
\end{equation}
which measures the angular power in multipoles above
$\ell_0$.  For the quadrupole and above ($\ell_0 = 2$):
$d_{\geq 2} < 2 \times 10^{-3}$ at all dipole-anomaly
scenarios.  The $\Teff$ manifold remains tangent to the exact
Boltzmann solution within this bound, and the representation
error is quadratic in $d_{\geq 2}$---of order $10^{-6}$.


% ============================================================
\section{Vorticity extension}\label{sec:teff-vort}
% ============================================================

\subsection{Azimuthal orthogonality}\label{sec:azimuth}

The $m = 1$ modes generated by vorticity ($W$ and $V$ in
Eq.~\eqref{eq:Theta-expansion}) are orthogonal to the $m = 0$
modes of the shear-sourced pattern under the angular
integration.  For any integer $n$:
\begin{equation}\label{eq:azimuth-orth}
\int_0^{2\pi}\cos^n\!\phi\,d\phi = 0
\quad\text{for $n$ odd}\,.
\end{equation}
All cross-terms between the axisymmetric ($m = 0$) shear
pattern and the non-axisymmetric ($m = 1$) vorticity pattern
vanish exactly in $\langle\Theta^\alpha\rangle_\Omega$ for any
power $\alpha$.  The vorticity enters only through the additive
variance terms $W^2/3$ and $V^2/15$ in
$\Sigma_2$~\eqref{eq:Sigma2-def}.


\subsection{Magnitude of the vorticity correction}
\label{sec:vort-magnitude}

In a spatially homogeneous barotropic model, the matter-frame
vorticity satisfies
$|\hat{\omega}|/\Theta \lesssim (\sigma/\Theta)\,\beta$
(Section~\ref{sec:king-ellis}), making it the product of two
small quantities.  The vorticity contribution to the moment map
is therefore suppressed by four powers of the fundamental
smallness parameter:
\begin{equation}\label{eq:vort-suppression}
\frac{\Delta\langle\Theta^4\rangle^{(\omega)}}
     {\langle\Theta^4\rangle^{(\sigma)} - 1}
\lesssim
\left(\frac{|\hat{\omega}|/\Theta}{\sigma/\Theta}\right)^2
\sim \beta^2
\sim 10^{-6}\,.
\end{equation}
At the S3 scenario, the absolute vorticity correction is
$\lesssim 3 \times 10^{-10}$---ten orders of magnitude below
the shear-induced departure.

The Gaunt coefficient $g_{224}$ is unmodified by vorticity (the
$A^2 Q$ cross-term involves only $m = 0$ modes), and the
tangency diagnostic changes by less than $10^{-10}$.  The
existing $\Teff$ formalism therefore requires no modification
for the inclusion of vorticity as an independent source.


% ============================================================
\section{Acceleration extension}\label{sec:teff-accel}
% ============================================================

\subsection{The monopole--dipole coupling}\label{sec:mono-dip}

Four-acceleration enters the $\Teff$ framework through the
$\ell = 0 \leftrightarrow \ell = 1$ coupling in the Boltzmann
hierarchy (Eq.~\eqref{eq:ell1}).  The source term
$\frac{4}{3}I\,\dot{u}_a\,\delta_{\ell 1}$ couples the
radiation monopole to the dipole with coefficient $4/3$,
generating a dipole anisotropy of amplitude
$\eps_1^{(\dot{u})} = \etaud\,\beta$
(Section~\ref{sec:accel-dipole}).

Unlike vorticity, the acceleration contribution is aligned with
the shear ($m = 0$) and adds coherently to the $A$ parameter in
the moment map:
\begin{equation}\label{eq:A-eff}
A_{\mathrm{eff}} = A_{\mathrm{tilt}} + A_{\dot{u}}
= \beta + \etaud\,\beta = (1 + \etaud)\,\beta\,.
\end{equation}
The acceleration increases the effective dipole by
$\sim 8.3\%$, a correction of the same order as the nonlinear
$\Teff$ correction itself.


\subsection{Cross-term hierarchy}\label{sec:cross-terms}

The acceleration--shear cross-term in $\langle\Theta^4\rangle$
contributes through the modified $A_{\mathrm{eff}}^2 Q$ coupling:
\begin{equation}
\delta\langle\Theta^4\rangle^{(\dot{u})}
= \tfrac{12}{5}\,g_{224}\,
  \bigl(A_{\mathrm{eff}}^2 - A_{\mathrm{tilt}}^2\bigr)\,Q
= \tfrac{12}{5}\,g_{224}\,
  \bigl(2\,\etaud + \etaud^2\bigr)\,\beta^2\,Q\,.
\end{equation}
The fractional contribution relative to the shear correction is
\begin{equation}
\frac{\delta R^{(\dot{u})}}{\Delta_\sigma}
\approx 2\,\etaud = 2/12 \approx 0.167
\quad (w = \tfrac{1}{3})\,,
\end{equation}
at leading order in $\etaud$.  In the definition of the full
correction factor, this translates to
$\delta R^{(\dot{u})} \approx 0.077\,\Delta_\sigma$
(approximately 7.7\% of the shear-only correction), confirmed by
direct numerical quadrature in the VN-04 computation.

The acceleration--vorticity cross-term is further
suppressed by two powers of $\beta$ (one from each) and is
negligible ($< 10^{-15}$ at S3).


% ============================================================
\section{Nonlinear $\ell$-pole mixing}\label{sec:ell-mixing}
% ============================================================

\subsection{The boost mixing matrix}\label{sec:boost-mix}

A Lorentz boost of rapidity $\beta$ along the symmetry axis
transforms the intrinsic quadrupole $Q_{\mathrm{int}}$ into a
superposition of all multipoles.  The leading-order mixing is
\begin{equation}\label{eq:ell-mixing}
\delta a_1^{(\mathrm{boost})}
= -\tfrac{4}{5}\,\beta\,Q_{\mathrm{int}}\,,\qquad
\delta a_3^{(\mathrm{boost})}
= +\tfrac{9}{5}\,\beta\,Q_{\mathrm{int}}\,,
\end{equation}
with second-order corrections
$\delta a_0, \delta a_2, \delta a_4 \propto \beta^2
Q_{\mathrm{int}}$.  The dipole leakage
$\delta a_1 \propto \beta Q$ is the mechanism by which shear
generates an apparent dipole through the observer's tilt, and
the octupole leakage $\delta a_3$ redistributes quadrupole
power to higher multipoles.

At the CF4 rapidity $\beta = 1.334 \times 10^{-3}$, the exact
and perturbative $O(\beta)$ mixing coefficients agree to within
$6\%$ for $\ell = 1$ and $1\%$ for $\ell = 3$.  The
perturbative treatment is adequate for the present analysis.


\subsection{Vorticity and acceleration mixing}
\label{sec:vort-accel-mix}

Vorticity generates $\Delta\ell = \pm 1$ mixing through
odd-parity ($m = 1$) modes.  The coupling involves the
Levi-Civita structure $\varepsilon_{bc\langle a_\ell}
\omega^b I^c{}_{A_{\ell-1}\rangle}$, which connects $\ell$ to
$\ell \pm 1$ with an alternating sign under parity.  The
physical origin is frame dragging: vorticity rotates the photon
direction, generating azimuthal structure in the temperature
field.

Acceleration generates $\Delta\ell = \pm 1$ mixing through
even-parity ($m = 0$) modes via the coupling
$\dot{u}_{\langle a_\ell} I_{A_{\ell-1}\rangle}$.  The parity
distinction has a decisive consequence: vorticity-induced modes
are orthogonal to the shear pattern and do not interfere with
the $\Rsig$ correction, while acceleration-induced modes add
coherently to the shear correction (both are axisymmetric).

The complete hierarchy of corrections to the moment map, ordered
by magnitude at the S3 scenario:
\begin{center}\footnotesize
\begin{tabular}{llcc}
\toprule
Source & Coupling & Magnitude
  & Ratio to $\langle\Theta^4\rangle - 1$ \\
\midrule
Shear angular variance & $\sigma^2/\Theta^2$ &
  $\sim 2 \times 10^{-5}$ & 1 \\
Shear self-coupling ($g_{224}$) & $A^2 Q$ &
  $\sim 10^{-9}$ & $\sim 10^{-5}$ \\
Accel.\ cross-term & $A_{\dot{u}}^2 Q$ &
  $\sim 7.7\%\,\Delta_\sigma$ & $\sim 10^{-6}$ \\
Vorticity dipole ($W^2$) & $\omega^2/\Theta^2$ &
  $\lesssim 10^{-15}$ & $\lesssim 10^{-10}$ \\
\bottomrule
\end{tabular}
\end{center}


% ============================================================
\section{The Einstein--$\Teff$ ODE system}\label{sec:teff-ode}\label{sec:transfer}
% ============================================================

\subsection{Extended dynamical system}\label{sec:ode-system}

The coupled evolution of shear, vorticity, tilt, and
deceleration parameter in the 1+3 covariant framework, expressed
in Hubble-normalised variables, is governed by the system
\begin{align}
\dot{\sigma}_H &= -(2 - q)\,H\,\sigma_H + S_\sigma\,,
\label{eq:sigma-ode}\\
\dot{\omega}_H &= -(2 - \tfrac{3}{2}w)\,H\,\omega_H + S_\omega\,,
\label{eq:omega-ode}\\
\dot{\beta} &= -\tfrac{1}{3}(1 - 3w)\,H\,\beta + S_\beta\,,
\label{eq:beta-ode}
\end{align}
where $\sigma_H = |\sigma|/\Theta$,
$\omega_H = |\omega|/\Theta$, and $S_\sigma$, $S_\omega$,
$S_\beta$ are source terms from the $\Teff$ back-reaction and
the Einstein constraint equations.  The shear source
$S_\sigma$ includes the Weyl tensor $E_{ab}$ (tidal field) and
the nonlinear $\Teff$ correction; the vorticity source
$S_\omega$ includes the shear--vorticity coupling
$\sigma_{c[a}\omega^c{}_{b]}$; the tilt source $S_\beta$
includes the acceleration feedback.


\subsection{Radiation-era behaviour}\label{sec:rad-era}

In the radiation era ($w = 1/3$, $q = 1$), the decay rates are:
\begin{equation}\label{eq:rad-decay}
\sigma_H \propto a^{-(2-q)} = a^{-1}\,,\qquad
\omega_H \propto a^{-(2-1/2)} = a^{-3/2}\,,\qquad
\dot{\beta} = 0\,.
\end{equation}
The tilt rapidity is \emph{frozen} during the radiation era---a
structural consequence of the equation of state $w = 1/3$
setting the $(1 - 3w)$ prefactor to zero.  In the standardised
variables, $\Sigstd$ is constant during radiation domination
(both numerator and denominator of
$\sigma_{ab}\sigma^{ab}/(6H^2)$ scale as $a^{-4}$), while
$\Wstd$ decays as $a^{-3}$.


\subsection{The quadrupole evolution equation}
\label{sec:Q-dot}

The effective quadrupole $Q$ evolves according to the $\ell = 2$
MDE~\eqref{eq:ell2}, which in the free-streaming regime
(post-decoupling) gives
\begin{equation}\label{eq:Q-ode}
\dot{Q} + \tfrac{4}{3}\Theta\,Q
= \tfrac{8}{15}\,\sigma_H\,\Theta
  - \tilde{D}^c I_{c\langle ab\rangle}/I
  + \tfrac{2}{5}\tilde{D}_{\langle a}I_{b\rangle}/I\,.
\end{equation}
For spatially homogeneous models, the gradient terms vanish on
the group orbits, and the quadrupole is driven solely by the
shear.  The Einstein--$\Teff$ ODE system couples this equation
to the shear evolution~\eqref{eq:sigma-ode} through the Weyl
tensor $E_{ab}$, which in turn depends on $Q$ through the
Boltzmann hierarchy.


% ============================================================
\section{Comprehensive nonlinear correction}
\label{sec:comprehensive}
% ============================================================

\subsection{The correction factors $\Rsig$, $\Rom$, $\Rud$}
\label{sec:R-factors}

The nonlinear correction to the MES shear bound is defined as
the ratio of the nonlinear to linearised bound:
\begin{equation}\label{eq:Rsig-def}
\Rsig := \frac{\Bsig^{\mathrm{nl}}}{\Bsig^{\mathrm{lin}}}
= 1 + \Delta_\sigma\,,
\end{equation}
where $\Delta_\sigma$ is the fractional correction.  Analogous
definitions hold for $\Rom = 1 + \Delta_\omega$ and
$\Rud = 1 + \Delta_{\dot{u}}$.

The $\Rsig$ correction is computed by evaluating the $\Theta^4$
moment map~\eqref{eq:F4-full} with the full nonlinear
temperature field and projecting onto the $\eps_\ell$ basis to
extract the effective bound.  The computation uses
Gauss--Legendre quadrature with 400 nodes, achieving machine
precision for the angular integrals.


\subsection{Universal ratios}\label{sec:universal}

The most remarkable feature of the nonlinear corrections is the
near-exact universality of the ratios between the three
correction factors:
\begin{equation}\label{eq:universal-ratios}
\frac{\Delta_\omega}{\Delta_\sigma}
= 1.478 \pm 0.002\,,\qquad
\frac{\Delta_{\dot{u}}}{\Delta_\sigma}
= 0.742 \pm 0.001\,,
\end{equation}
where the uncertainties reflect the scenario dependence across
S1--S3.  These ratios are not fitted parameters but emerge from
the structure of the $\Theta^4$ moment map: the Clebsch--Gordan
coefficients that couple different $\ell$-modes determine the
relative sensitivity of each bound to the nonlinear corrections,
and these coefficients are scenario-independent.

The physical content of the universality is that all three
corrections are driven by the same underlying mechanism---the
$A^2 Q$ coupling in $\langle\Theta^4\rangle$---and differ only
in the projection coefficients that relate this coupling to the
individual bound combinations.


\subsection{Correction table}\label{sec:correction-table}

Table~\ref{tab:nl-corrections} collects the nonlinear
corrections for all scenarios.  The key features are:
$\Rsig(\mathrm{S0}) = 1$ exactly (no anisotropy, no
correction); $\Rsig$ increases monotonically with $\eps_1$;
the fractional correction to $\Sigstd$ reaches $0.8\%$ at S3
($\eps_1 = 1.476 \times 10^{-3}$) and $1.8\%$ at the S2c
radio-dipole amplitude ($\eps_1 = 3.296 \times 10^{-3}$);
and the hierarchy $\Rom > \Rsig > \Rud$ is preserved across all
scenarios.

\begin{table}[h]
\caption{Nonlinear $\Teff$ corrections by scenario.  $\Rsig$,
$\Rom$, $\Rud$ are the correction factors for shear, vorticity,
and acceleration bounds; $\delta\Sigstd/\Sigstd = \Rsig^2 - 1$
is the fractional correction to the standardised shear.}
\label{tab:nl-corrections}
\centering\footnotesize\setlength{\tabcolsep}{4pt}
\begin{tabular}{lcccccc}
\toprule
Scen.\ & $\eps_1$ & $\Rsig$ & $\Rom$ & $\Rud$
  & $\delta\Sigstd/\Sigstd$ (\%) \\
\midrule
S0  & $0$             & 1.000000 & 1.000000 & 1.000000 & 0.000 \\
S1  & $1.233\times10^{-3}$ & 1.003316 & 1.004898 & 1.002460 & 0.664 \\
S2a & $1.476\times10^{-3}$ & 1.003971 & 1.005866 & 1.002947 & 0.796 \\
S2c & $3.296\times10^{-3}$ & 1.008893 & 1.013155 & 1.006595 & 1.786 \\
\bottomrule
\end{tabular}
\end{table}


% ============================================================
\section{Defect propagation}\label{sec:defect-prop}
% ============================================================

\subsection{Shear channel}\label{sec:shear-channel}

The nonlinear correction to $\Sigstd$ is
$\delta\Sigstd/\Sigstd = \Rsig^2 - 1$, reaching $0.80\%$ at
S3 ($\eps_1 = 1.476 \times 10^{-3}$).  The correction propagates directly to $\xI = \Sigstd$ for
type~I and to $\xmax$ for all types through the MES ceiling:
$x^{\mathrm{max,nl}} = (3/2)(\Rsig\,\Bsig)^2
= \Rsig^2\,x^{\mathrm{max,lin}}$.


\subsection{Tilt channel}\label{sec:tilt-channel}

The nonlinear correction modifies the effective dipole amplitude
$\eps_1^{\mathrm{eff}}$, which changes $\beta$ through the
safe-route relation.  The correction is
$\delta\beta/\beta \sim 10^{-6}$ at S1---four orders of
magnitude smaller than the frame-attribution correction
($\sim 8\%$).  The tilt correction is therefore negligible
compared to the VT-07 frame correction.


\subsection{Combined assessment}\label{sec:combined}

Table~\ref{tab:defect-prop} shows the defect bounds with and
without nonlinear corrections for types~I and~V.

\begin{table}[h]
\caption{Defect bounds with nonlinear corrections.
$\delta x/x$ is the fractional change from the linearised bound.
$\xV/\xI$ is the tilt fraction.}
\label{tab:defect-prop}
\centering\footnotesize\setlength{\tabcolsep}{3.5pt}
\begin{tabular}{lcccccc}
\toprule
Scen.\ & $\xI^{\mathrm{lin}}$ & $\xI^{\mathrm{nl}}$
  & $\delta\xI/\xI$ (\%)
  & $\xV^{\mathrm{lin}}$ ($w{=}0$)
  & $\xV^{\mathrm{nl}}$ ($w{=}0$)
  & $\xV/\xI$ \\
\midrule
S0  & $2.64\times10^{-10}$ & $2.64\times10^{-10}$ & 0.000
    & 0 & 0 & --- \\
S1  & $6.42\times10^{-6}$  & $6.46\times10^{-6}$  & 0.664
    & $4.09\times10^{-7}$ & $4.09\times10^{-7}$ & 0.063 \\
S3  & $9.18\times10^{-6}$  & $9.25\times10^{-6}$  & 0.796
    & $5.86\times10^{-7}$ & $5.86\times10^{-7}$ & 0.063 \\
\bottomrule
\end{tabular}
\end{table}

The nonlinear correction to $\xI$ is $\leq 0.8\%$ across
S0--S3, while the correction to $\xV$ (the tilt
contribution) is $\leq 0.002\%$---dominated by the shear
channel.

The small-$\beta$ approximation
$\sinh^2\!\beta \approx \beta^2$ used throughout the analysis
is validated by the exact computation.  At the S3 scenario
($\beta = 1.36 \times 10^{-3}$), the fractional error
$(\sinh^2\!\beta - \beta^2)/\beta^2 \approx 2 \times 10^{-6}$
is five orders of magnitude below the nonlinear $\Teff$
correction and twelve orders below the leading defect bound.


\input{chapters/figures_ch05}
